\documentclass[12pt,a4paper]{article}

\usepackage[T1]{fontenc}
\usepackage[utf8]{inputenc}
\usepackage[top=2.5cm, bottom=2.5cm, left=3cm, right=3cm]{geometry}
\usepackage{microtype}
\usepackage{parskip}
\usepackage{enumitem}

\pagestyle{empty}

\begin{document}

% ---------------------------------------------------------------
% Sender / contact
% ---------------------------------------------------------------
\begin{flushright}
David Andersson\\
Chalmers Next Labs\\
Chalmers University of Technology\\
Gothenburg, Sweden\\[0.4em]
\texttt{david.andersson@chalmers.se}\\[1em]
\today
\end{flushright}

\bigskip

% ---------------------------------------------------------------
% Body
% ---------------------------------------------------------------
Dear Editors,

We are pleased to submit our manuscript, ``Low-carbon lifestyles show substantially lower emissions with no sign of rebound,'' for consideration as an Article in \textit{Communications Sustainability}. The manuscript is transferred from \textit{Nature Climate Change} (NCLIM-26093896), whose editors kindly arranged for external review at your journal. Please note that the files accompanying this transfer have been substantially revised since the original submission; this letter describes the current version.

How much can lifestyle change contribute to sustainable consumption and climate mitigation? The most recent global assessment puts the direct potential at roughly 10~GtCO$_2$e per year among high-consuming households, but assumes that the money saved is re-spent on other consumption, erasing around half of this potential. This re-spending assumption, inherited from energy-efficiency analysis, has become a standard default in demand-side mitigation modelling but has never been tested against observed consumption data. We provide that test.

Linking twelve months of bank-transaction records---covering close to the full consumption budget---with administrative registers and survey responses for a random sample of 1,297 Swedish single-adult households, we observe, rather than simulate, the complete spending and emission patterns of established car-free, flight-free, and meat-free households, comparing them with demographically similar non-adopters.

We believe three findings will interest your readership. First, lifestyle adopters show large net emission reductions---adjusted total differences of $-$1.51 and $-$2.06~tCO$_2$e per person-year ($-$25\% and $-$26\%) for car-free and flight-free living. Second, there is no sign of automatic rebound: no lifestyle group shows higher indirect emissions, and for both mobility lifestyles the observed indirect emissions fall significantly below proportional re-spending benchmarks ($p = 0.041$ and $p = 0.045$). Third, motivation shapes the outcome: households with stronger environmental self-identity show lower emissions across the entire budget, so that environmentally motivated adopters achieve the broadest reductions.

These results speak directly to the design of demand-side sustainability policy: they indicate that the mitigation potential of low-carbon lifestyles is not inherently self-eroding, and that the fate of freed-up income depends on who adopts and why. The study combines transaction-based carbon accounting, registry data, and environmental psychology, and we hope its interdisciplinary approach will resonate with the breadth of \textit{Communications Sustainability}'s readership.

The work is original and not under consideration elsewhere. We declare one competing interest, detailed in the manuscript: David Andersson is also the founder of Svalna AB, which developed the mobile app and data-collection infrastructure used in the study (and has since ceased consumer-facing operations). All statistical analysis, interpretation of results, and manuscript preparation were conducted independently of Svalna AB.

% ---------------------------------------------------------------
% Suggested reviewers
% ---------------------------------------------------------------
We would respectfully suggest the following potential reviewers:

\begin{itemize}[noitemsep, topsep=0.5em, leftmargin=1.5em]
  \item Felix Creutzig --- MCC Berlin / TU Berlin, Germany
  \item Diana Ivanova --- Sustainability Research Institute,
        University of Leeds, UK
  \item Steve Sorrell --- Sussex Energy Group / SPRU,
        University of Sussex, UK
  \item John Th{\o}gersen --- Department of Management,
        Aarhus University, Denmark
  \item Lorraine Whitmarsh --- Department of Psychology,
        University of Bath, UK
\end{itemize}

Thank you for considering our work.

\bigskip

Sincerely,

David Andersson\\[0.2em]
\textit{Corresponding author, on behalf of all co-authors}\\[0.2em]
Chalmers Next Labs, Chalmers University of Technology, Gothenburg, Sweden\\[0.2em]
\texttt{david.andersson@chalmers.se}

\end{document}
